\begin{textsample}{POS Dim 3 – persona\_gemini – Score 19.00 – t572\_gemini.txt}  \label{ex:f3_pos_042}
A culture of extravagant \textbf{consumption} is creating a massive environmental \textbf{problem}, with annual electronic \textbf{waste} now exceeding 65 million \textbf{tonnes}.
This crisis is driven by a \textbf{system} of planned obsolescence, where \textbf{companies} like Apple \textbf{design} \textbf{products}, such as iPhones, that are intentionally hard to repair.
These corporations also actively lobby against right-to-repair laws, ensuring \textbf{consumers} are forced into a \textbf{cycle} of frequent replacement.
The life \textbf{cycle} of these devices begins with the destructive mining of \textbf{minerals} and \textbf{ends} with a tragically short lifespan, contributing directly to the growing e-waste mountain.
A Greenpeace \textbf{report} has graded electronics firms on their commitment to sustainability.
While \textbf{companies} such as HP, Dell, and Fairphone received praise for improving the repairability of their \textbf{products}, others, including Samsung, were criticized for their poor performance.
To address this, we must \textbf{shift} from our current \textbf{model} to one of circular \textbf{production}.
This change \textbf{needs} to be reinforced by government policies that mandate durability and repairability.
Embracing more modest \textbf{consumption} habits is also crucial to reduce our ecological overshoot and support the vital \textbf{growth} of renewable \textbf{energy}.

% matched lemmas: company, consumer, consumption, cycle, design, end, energy, growth, mineral, model, need, problem, product, production, report, shift, system, tonne, waste
\end{textsample}
